\documentclass[9pt,a4paper,twocolumn]{ctexart}

% --- Packages ---
\usepackage[top=1.2cm, bottom=1.2cm, left=1.2cm, right=1.2cm, columnsep=0.8cm]{geometry}
\usepackage{hyperref}
\usepackage{graphicx}
\usepackage{float}
\usepackage{longtable}
\usepackage{tabularx}
\usepackage{booktabs}
\usepackage{enumitem}
\usepackage{xcolor}
\usepackage{fancyhdr}
\usepackage{listings}
\usepackage{fontspec}
\usepackage{titlesec}
\usepackage{caption}

% --- Code block styling ---
\definecolor{codebg}{RGB}{245,245,245}
\definecolor{codeframe}{RGB}{220,220,220}
\definecolor{codekw}{RGB}{0,0,192}
\definecolor{codecomment}{RGB}{0,128,0}
\definecolor{codestring}{RGB}{163,21,21}
\definecolor{codenum}{RGB}{128,0,128}
\definecolor{codepunct}{RGB}{90,90,90}
\definecolor{badgegreen}{RGB}{34,139,34}
\definecolor{badgeblue}{RGB}{30,144,255}
\definecolor{badgeblack}{RGB}{50,50,50}

% --- Difficulty badges (rounded corners via tikz, pre-rendered in saveboxes) ---
\usepackage{tikz}
\newsavebox{\badgechubox}
\savebox{\badgechubox}{\tikz[baseline]{\node[fill=badgegreen,text=white,rounded corners=2.5pt,inner xsep=4pt,inner ysep=1.5pt,font=\sffamily\footnotesize\bfseries]{初级};}}
\newsavebox{\badgezhongbox}
\savebox{\badgezhongbox}{\tikz[baseline]{\node[fill=badgeblue,text=white,rounded corners=2.5pt,inner xsep=4pt,inner ysep=1.5pt,font=\sffamily\footnotesize\bfseries]{中级};}}
\newsavebox{\badgegaobox}
\savebox{\badgegaobox}{\tikz[baseline]{\node[fill=badgeblack,text=white,rounded corners=2.5pt,inner xsep=4pt,inner ysep=1.5pt,font=\sffamily\footnotesize\bfseries]{高级};}}
\newcommand{\lvchu}{\usebox{\badgechubox}}
\newcommand{\lvzhong}{\usebox{\badgezhongbox}}
\newcommand{\lvgao}{\usebox{\badgegaobox}}
\pdfstringdefDisableCommands{%
  \def\lvchu{[初级]}%
  \def\lvzhong{[中级]}%
  \def\lvgao{[高级]}%
}

\lstset{
  basicstyle=\ttfamily\scriptsize,
  keywordstyle=\color{codekw}\bfseries,
  commentstyle=\color{codecomment}\itshape,
  stringstyle=\color{codestring},
  backgroundcolor=\color{codebg},
  frame=single,
  rulecolor=\color{codeframe},
  breaklines=true,
  breakatwhitespace=false,
  breakindent=0pt,
  postbreak=\mbox{\textcolor{red}{$\hookrightarrow$}\space},
  numbers=left,
  numberstyle=\tiny\color{gray},
  columns=flexible,
  keepspaces=true,
  showstringspaces=false,
  tabsize=2,
  aboveskip=4pt,
  belowskip=4pt,
  extendedchars=true,
  literate={，}{{，}}1 {；}{{；}}1 {！}{{！}}1 {？}{{？}}1 {“}{{“}}1 {”}{{”}}1 {（}{{（}}1 {）}{{）}}1 {《}{{《}}1 {》}{{》}}1,
}

% --- Extra language definitions (no built-in listings support) ---
\lstdefinelanguage{json}{
  morekeywords={true,false,null},
  sensitive=true,
  morestring=[b]",
  comment=[l]{//},
  morecomment=[s]{/*}{*/},
}
\lstdefinelanguage{yaml}{
  morekeywords={true,false,null,yes,no,on,off},
  sensitive=false,
  morecomment=[l]{\#},
  morestring=[b]",
  morestring=[d]',
}
\lstdefinelanguage{http}{
  morekeywords={GET,POST,PUT,DELETE,PATCH,HEAD,OPTIONS,HTTP,Host,Accept,Authorization,Content-Type,Content-Length,User-Agent},
  sensitive=true,
  morecomment=[l]{\#},
  morestring=[b]",
}
\lstdefinelanguage{TypeScript}{
  morekeywords={abstract,any,as,async,await,break,case,catch,class,const,continue,debugger,declare,default,delete,do,else,enum,export,extends,false,finally,for,from,function,get,if,implements,import,in,instanceof,interface,keyof,let,module,namespace,never,new,null,number,of,package,private,protected,public,readonly,return,set,static,string,super,switch,symbol,this,throw,true,try,type,typeof,undefined,unknown,var,void,while,with,yield,Record,Array,Promise,AsyncIterable,ReadableStream,Date,Error,Object,JSON},
  sensitive=true,
  morecomment=[l]{//},
  morecomment=[s]{/*}{*/},
  morestring=[b]",
  morestring=[b]',
  morestring=[b]`{`},
}

% --- Compact spacing ---
\linespread{0.95}
\setlength{\parskip}{1pt plus 1pt minus 1pt}
\setlength{\parindent}{0pt}

% --- Table styling ---
\renewcommand{\arraystretch}{1.1}

% --- Heading styling (numbered) ---
\titleformat{\section}{\normalsize\bfseries}{\thesection\quad}{0em}{}
\titlespacing{\section}{0pt}{6pt}{2pt}
\titleformat{\subsection}{\small\bfseries}{\thesubsection\quad}{0em}{}
\titlespacing{\subsection}{0pt}{4pt}{1pt}
\titleformat{\subsubsection}{\small\bfseries}{\thesubsubsection\quad}{0em}{}
\titlespacing{\subsubsection}{0pt}{3pt}{1pt}

% --- Page style ---
\pagestyle{fancy}
\fancyhf{}
\fancyhead[L]{\small\emph{\leftmark}}
\fancyhead[R]{\small\thepage}
\renewcommand{\headrulewidth}{0.4pt}
\renewcommand{\sectionmark}[1]{}  % prevent sections from overwriting leftmark

% --- Hyperref setup ---
\hypersetup{
  colorlinks=true,
  linkcolor=blue,
  urlcolor=blue,
  citecolor=blue,
}

% --- Caption format ---
\DeclareCaptionFormat{myformat}{\small\bfseries #1#2#3}
\captionsetup{format=myformat}

\begin{document}

\title{多智能体编排：模式、通信与失败边界}
\date{}
\maketitle
\markboth{TypeScript 版 Agent 知识}{ TypeScript 版 Agent 知识}

\section*{适用条件}

\begin{itemize}[itemsep=1pt, topsep=2pt]
  \item 子任务需要不同的 System Prompt、工具集或记忆边界；
  \item 子任务可并行，且互不阻塞；
  \item 需要生成、审查、验证角色分离；
  \item 不同角色由不同团队维护，需要独立发布与权限隔离。
\end{itemize}


先用单 Agent + Workflow 验证；满足以上条件后再上多 Agent。



\section*{编排模式}


\subsection*{1. Orchestrator-Workers（编排者-工作者）}


\textbf{结构}：一个 Orchestrator 负责拆解任务、分发给多个 Worker，Worker 返回结构化结果，Orchestrator 汇总并交付。



\begin{lstlisting}
用户任务
   │
   ▼
Orchestrator ──任务A──▶ Worker A
   │  │
   │  └──────任务B──▶ Worker B
   │  │
   │  └──────任务C──▶ Worker C
   │
   ▼
汇总结果
\end{lstlisting}


\begin{itemize}[itemsep=1pt, topsep=2pt]
  \item 优点：并行度高，子任务互不干扰，扩展容易。
  \item 缺点：Orchestrator 是单点，拆解质量决定全局上限。
  \item 适合：调研、报告生成、多源信息汇总、多模块代码生成。
  \item Worker 通常设计为无状态：输入任务卡，输出结果卡，不自行决定全局。
\end{itemize}


\subsection*{2. Supervisor（监督者）}


\textbf{结构}：Supervisor 不拆全局计划，而是根据当前状态，在多个专业 Agent 之间路由，直到满足退出条件。



\begin{lstlisting}
          Supervisor
          │    │    │
          ▼    ▼    ▼
      Agent A Agent B Agent C
          │    │    │
          └────┴────┘
            返回结果
\end{lstlisting}


\begin{itemize}[itemsep=1pt, topsep=2pt]
  \item 优点：动态路由，能根据中间结果换角色；比 Orchestrator 更灵活。
  \item 缺点：监督逻辑容易变成难以测试的 Prompt 黑洞。
  \item 适合：诊断、客服升级、复杂问题需要多轮角色切换的场景。
  \item 与 Orchestrator 的区别：Orchestrator 先拆后发，Supervisor 边看边路由。
\end{itemize}


\subsection*{3. Swarm / Handoff（动态交接）}


\textbf{结构}：没有中心调度者。Agent 之间按规则或模型判断进行交接，一次只有一个 Agent 活跃，直到任务完成。



\begin{lstlisting}
Agent A ──交接──▶ Agent B ──交接──▶ Agent C
  ▲                                        │
  └────────────── 任务完成 ◀───────────────┘
\end{lstlisting}


\begin{itemize}[itemsep=1pt, topsep=2pt]
  \item 优点：上下文随交接流转，交互自然，适合对话式客服。
  \item 缺点：交接判定难调，容易无限交接或上下文膨胀。
  \item 适合：客服转接、多领域助手、需要“把当前用户带到正确专家”的场景。
  \item 必须设计：最大交接次数、交接时携带的最小上下文、无路可走时的兜底 Agent。
\end{itemize}


\subsection*{4. Hierarchical（分层递归）}


\textbf{结构}：高层 Agent 拆解任务后，子 Agent 可以继续拆解并生成下一级 Agent，形成树形递归结构。



\begin{lstlisting}
                顶层 Agent
          ┌─────────┼─────────┐
          ▼         ▼         ▼
      子 Agent   子 Agent   子 Agent
                    │
              ┌─────┴─────┐
              ▼           ▼
         孙 Agent     孙 Agent
\end{lstlisting}


\begin{itemize}[itemsep=1pt, topsep=2pt]
  \item 优点：适合大而深的复杂任务，如大型代码库重构、多系统迁移方案。
  \item 缺点：成本与延迟随层数指数上升，递归终止条件极难控制。
  \item 规则：\textbf{层数默认不超过 2 层}；每一层必须有预算和期限。
\end{itemize}


\subsection*{5. Blackboard（黑板模式）}


\textbf{结构}：多个专业 Agent 共享一个结构化工作区（黑板），各自读取相关部分、写入结论，由控制组件或共识规则收敛结果。



\begin{lstlisting}
Agent A ─写─▶ ┌──────────────┐ ◀─读─ Agent C
              │   Blackboard  │
Agent B ─写─▶ │  共享状态/事实 │ ◀─读─ Agent D
              └──────────────┘
\end{lstlisting}


\begin{itemize}[itemsep=1pt, topsep=2pt]
  \item 优点：解耦彻底，适合需要多学科知识融合的开放问题。
  \item 缺点：黑板内容管理难，冲突消解复杂，最容易被忽视。
  \item 适合：诊断、研究综述、方案评审。
  \item 关键设计：写入权限、事实/假设分区、冲突解决规则、收敛截止时间。
\end{itemize}


\subsection*{6. Peer-to-Peer（对等协作）}


\textbf{结构}：没有中心调度者，Agent 之间平等对话、互相审查或协商。


\begin{itemize}[itemsep=1pt, topsep=2pt]
  \item 优点：能产生对抗式质量提升，适合评审、辩论、模拟谈判。
  \item 缺点：通信开销最大，最难观测与终止。
  \item 使用建议：只作为局部审查结构，不让它成为全局拓扑。
\end{itemize}


\section*{模式对比表}

{\footnotesize
\begin{tabular}{|p{\dimexpr 0.116\columnwidth-2\tabcolsep\relax}|p{\dimexpr 0.290\columnwidth-2\tabcolsep\relax}|p{\dimexpr 0.145\columnwidth-2\tabcolsep\relax}|p{\dimexpr 0.072\columnwidth-2\tabcolsep\relax}|p{\dimexpr 0.174\columnwidth-2\tabcolsep\relax}|p{\dimexpr 0.145\columnwidth-2\tabcolsep\relax}|p{\dimexpr 0.058\columnwidth-2\tabcolsep\relax}|}
\hline
\textbf{维度} & \textbf{Orchestrator-Workers} & \textbf{Supervisor} & \textbf{Swarm} & \textbf{Hierarchical} & \textbf{Blackboard} & \textbf{P2P} \\
\hline
控制方式 & 中心拆解 & 中心路由 & 链式交接 & 树形递归 & 共享空间 & 无中心 \\
\hline
并行度 & 高 & 中 & 低 & 中高 & 高 & 中 \\
\hline
动态性 & 低 & 中 & 高 & 中 & 高 & 最高 \\
\hline
调试难度 & 低 & 中 & 中高 & 高 & 高 & 最高 \\
\hline
Token 开销 & 中 & 中高 & 高 & 很高 & 高 & 最高 \\
\hline
最适合 & 可并行的信息汇总 & 多角色动态诊断 & 客服转接 & 大型分层任务 & 多学科融合 & 评审辩论 \\
\hline
\end{tabular}
}



\section*{通信设计}


\subsection*{1. 不要用自然语言“聊天”}


Agent 之间自由聊天会产生三个问题：信息丢失、格式漂移、Token 爆炸。工程上应把 Agent 间通信当作\textbf{微服务接口}来设计。



\subsection*{2. 三种通信载体}

{\footnotesize
\begin{tabular}{|p{\dimexpr 0.170\columnwidth-2\tabcolsep\relax}|p{\dimexpr 0.319\columnwidth-2\tabcolsep\relax}|p{\dimexpr 0.255\columnwidth-2\tabcolsep\relax}|p{\dimexpr 0.255\columnwidth-2\tabcolsep\relax}|}
\hline
\textbf{载体} & \textbf{适用场景} & \textbf{优点} & \textbf{缺点} \\
\hline
共享 State & 同一进程 / 同一 Graph & 零序列化成本，调试直接 & 耦合强，边界不清 \\
\hline
结构化消息 & 跨进程 / 跨 Agent & 可校验、可版本化、可审计 & 需要定义 Schema \\
\hline
自然语言摘要 & 需要人读的中间产物 & 可解释性好 & 不稳定，不宜作为系统契约 \\
\hline
\end{tabular}
}



\subsection*{3. 一个最小交接 Payload}


\begin{lstlisting}[language=json]
{
  "task_id": "T-20260816-001",
  "from": "triage-agent",
  "to": "billing-agent",
  "intent": "refund_request",
  "context": {
    "order_id": "O-88231",
    "user_language": "zh-CN",
    "attempts": 1
  },
  "acceptance_criteria": [
    "确认订单状态",
    "判断是否满足 7 天无理由退款",
    "输出结构化处理建议"
  ],
  "max_tokens": 4000,
  "deadline_ms": 60000
}
\end{lstlisting}



\subsection*{4. A2A 与 MCP 的分工}

\begin{itemize}[itemsep=1pt, topsep=2pt]
  \item \textbf{MCP}：解决 Agent 与外部工具/资源之间的接入问题；
  \item \textbf{A2A}：解决 Agent 与 Agent 之间的任务与状态交换问题；
  \item 两者不冲突：Worker 可以一边通过 MCP 使用工具，一边通过 A2A 契约接收任务和回传结果。
\end{itemize}


\subsection*{5. 通信设计检查单}

\begin{itemize}[itemsep=1pt, topsep=2pt]
  \item 每条消息是否有 Schema 和版本？
  \item 是否限制上下文随交接无限膨胀？
  \item 是否记录消息轨迹并支持回放？
  \item 是否有超时、重试、幂等和死信处理？
  \item 是否禁止 Agent 之间直接共享密钥或不可审计的隐式状态？
\end{itemize}


\section*{状态与上下文边界}


多 Agent 系统最容易犯的错误是“全局共享一个上下文”。



推荐三条边界：


\begin{enumerate}[itemsep=1pt, topsep=2pt]
  \item \textbf{全局共享最小事实}：任务 ID、目标、最终交付物位置；
  \item \textbf{局部上下文各管各的}：每个 Agent 只拿自己的 System Prompt、工具描述和必要任务卡；
  \item \textbf{交接只传结果与验收标准}，不传完整推理过程，除非排查问题。
\end{enumerate}


\section*{失败模式库}

{\footnotesize
\begin{tabular}{|p{\dimexpr 0.105\columnwidth-2\tabcolsep\relax}|p{\dimexpr 0.421\columnwidth-2\tabcolsep\relax}|p{\dimexpr 0.474\columnwidth-2\tabcolsep\relax}|}
\hline
\textbf{失败模式} & \textbf{表现} & \textbf{防御} \\
\hline
目标漂移 & 子 Agent 越做越偏离原始任务 & 任务卡写清验收标准；Orchestrator 做终检 \\
\hline
无限交接 & 客服场景 A→B→C→A 循环 & 最大交接次数、全局兜底策略 \\
\hline
上下文稀释 & 消息越来越长，关键信息丢失 & 交接时只传结构化摘要 \\
\hline
重复劳动 & 多个 Agent 同时做同一子任务 & 任务卡带排他与幂等 ID \\
\hline
成本爆炸 & 并行度失控或递归过深 & 全局 Token 预算、层数与并发上限 \\
\hline
权限聚合 & 每个 Agent 权限都不大，组合后过大 & 全局最小权限，跨 Agent 动作单独审批 \\
\hline
单点失效 & Orchestrator 输出格式坏掉，全员空转 & Schema 校验、重试、降级为固定 Workflow \\
\hline
结果互相矛盾 & Blackboard/P2P 无法收敛 & 冲突消解规则、表决或仲裁 Agent、截止时间 \\
\hline
\end{tabular}
}



\section*{最小实现示例：Orchestrator-Workers}


\begin{lstlisting}
function run_orchestrator(task):
    plan = llm_plan(task, max_steps=5)          # 1. 生成任务卡列表
    cards = validate_plan_schema(plan)          # 2. Schema 校验，失败重试一次
    budget = Budget(max_total_tokens=50000)

    results = parallel_map(cards, worker_call)  # 3. 并行分发给 Worker
    report = llm_summarize(task, results)       # 4. 汇总结果
    return validate_report(report, task)        # 5. 对照验收标准终检

function worker_call(card):
    ctx = build_context(
        system=card.role_prompt,                # 只注入角色自己的提示词
        task=card.objective,
        acceptance=card.acceptance_criteria,
        tools=card.allowed_tools                # 只暴露白名单工具
    )
    for round in 1..card.max_rounds:
        action = llm_decide(ctx)                # 下一步：调工具或产出结果
        if action.is_final:
            return schema_validate(action.result, card.result_schema)
        observation = safe_execute(action.tool_call)
        ctx.append(observation)                 # 只追加观察，不追加闲聊
    raise WorkerTimeout(card)
\end{lstlisting}



这个骨架说明了多智能体系统与单 Agent 的关键差异：\textbf{计划、执行、校验被强制拆开，且每个 Worker 只能看到任务卡允许它看到的东西。}



\section*{选型决策表}

{\footnotesize
\begin{tabular}{|p{\dimexpr 0.245\columnwidth-2\tabcolsep\relax}|p{\dimexpr 0.377\columnwidth-2\tabcolsep\relax}|p{\dimexpr 0.377\columnwidth-2\tabcolsep\relax}|}
\hline
\textbf{场景特征} & \textbf{推荐模式} & \textbf{不推荐} \\
\hline
多个独立子任务可并行 & Orchestrator-Workers & P2P \\
\hline
客服多技能转接 & Swarm & Hierarchical \\
\hline
复杂诊断、角色动态切换 & Supervisor & Orchestrator-Workers \\
\hline
超大型分层任务 & Hierarchical（≤2 层） & Swarm \\
\hline
多学科融合、结果难预先定义 & Blackboard & Orchestrator-Workers \\
\hline
评审、辩论、对抗式质检 & P2P 局部使用 & 全局 P2P \\
\hline
流程可提前确定 & 不要多 Agent，用 Workflow & 任何多 Agent 模式 \\
\hline
\end{tabular}
}



\section*{落地十步清单}

\begin{enumerate}[itemsep=1pt, topsep=2pt]
  \item 用单 Agent + Workflow 跑通最小闭环；
  \item 从真实失败模式中找到必须拆分的边界；
  \item 选择\textbf{一种}中心化模式起步，默认 Orchestrator-Workers；
  \item 为任务卡和返回结果定义 Schema；
  \item 给每个 Agent 独立 System Prompt、工具白名单与上下文边界；
  \item 设置全局并发、层数、轮次、Token、超时预算；
  \item 记录消息轨迹，支持按 task\_id 回放；
  \item 先离线用评测集验证，再做影子流量；
  \item 准备好降级路径：多 Agent 失败回到单 Agent 或固定 Workflow；
  \item 每次新增 Agent 都要重新评估成本、权限和调试复杂度。
\end{enumerate}


\section*{相关文档}

\begin{itemize}[itemsep=1pt, topsep=2pt]
  \item \href{01-Agent核心概念.md}{AI Agent 核心概念}：Multi-Agent、A2A 基础概念
  \item \href{../04-工程实践/01-Workflow-Graph与Loop.md}{AI 工作流中的 Workflow、Graph 与 Loop}：什么时候用 Workflow
  \item \href{../04-工程实践/03-Harness工程.md}{一文搞懂 Harness Engineering}：Anthropic 三智能体架构案例
  \item \href{../00-概念与术语/01-Agent术语表与概念边界.md}{Agent 术语表与概念边界}：统一术语口径
\end{itemize}


\begin{quote}
一句话总结：\textbf{多智能体的价值来自边界清晰的分工，而不是 Agent 数量的堆叠。}
\end{quote}


\end{document}
